% Pursuing Stacks (cotes 134-2 to 134-8): archive cote 134-6, site page 31 (PDF page 32), typescript page 349.
% typescript. Draft compared: ps5.tex:2869-2916.
% AI-transcribed by Claude Opus 5.5 (claude-opus-5-5, adaptive thinking, effort high), Message Batches API, batch ps3 (msgbatch_01XdFmuVVBAcgc7EFxdtCqsj), 2026-11-12; not yet checked by eye.
\scanpage{134-6/31}
\mspage{349}

the condition that i be $M_a$-$\underline{W}$-aspheric is always satisfied !\hand{)} However,
\xout{whenxwe}\inserted[typed]{let's} assume A to be totally $\underline{W}$-aspheric and every object of A has a
section(over the final object $e_{A^{\wedge}}$ of $A^{\wedge}$ ) or, what amounts to the
same, every``non-empty''\hand{/}object of $A^{\wedge}$ has a section - \xout{thesexconditionsx}
\xout{inxfact} we'll say in this case A is ``\uline{strictly} \xout{Wxasph} \uline{totally $\underline{W}$-aspheric}\hand{''}
(compare section 60, p.149 , in the particular case $\underline{W} = \underline{W}_{\infty}$ , and
with $A^{\wedge}$ replaced by an arbitrary topos). Let's assume moreover that
$\underline{W}$ satisfies Loc 4) . In this case, the homotopy structure $h' = h_{\underline{W}_A}$
does not depend on the choice of $\underline{W}$ , namely it is the so-called
``canonical'' homotopy structure'' $h'' = h_{A^{\wedge}}$ of the \inserted[typed]{strictly} (totally 0-connected)
category $A^{\wedge}$ (cf. section 57), which in the special case of a category
$A^{\wedge}$ can also be defined as the homotopy structure $h_{\underline{W}_o\hand{A}}$ associated to
$\underline{W}_o\hand{A}$, where $\underline{W}_o$ is the coarsest basic localizer satisfying Loc 4) i.e.
\[ \underline{W}_o = \hand{\{} f \hand{\in} \mathrm{Fl}(\mathrm{Cat}) \hand{|} \hand{\pi}_o(f) \text{ bijective}\hand{\}} \quad . \]
The proof of this fact \xout{isxessentiallyxthexsamexas} $h' = h''$ i.e.
\[ h_{\underline{W}_A} = h_{A^{\wedge}} \quad ( = h_{\underline{W}_o} ) \]
is essentially the same as for the similiar prop. (section 60, p.149),
The condition Loc 4) on $\underline{W}$ i.e. $\underline{W} \hand{\subset} \underline{W}_o$ clearly implies
\[ h_{\underline{W}_A} \hand{\subset} h_{\underline{W}_o\hand{A}} \quad , \]
and to get the opposite inequality, for which we'll use the assumption
on A , we only have to prove that for any 0-connected object X of $A^{\wedge}$,
any two sections are $(h'=h_{\underline{W}_A}$ )-homotopic, a fortiori ( as $h \hand{\leq} h'$
by the assumption of total $\underline{W}$-asphericity of A) it is enough to prove
they are h-homotopic. Now this \deleted{follows from the following easy result}
\deleted{applied to \quad $= A^{\wedge}$ (and which implies equally the proposition of secti-}%
\deleted{on 60 cited above).} \inserted{\deleted{is stated} follows from lemma 2 p.268, applied
to $\mathcal{C} = A^{\wedge}$, $C = A$.}

\deleted{\uline{Proposition} 4. Let \quad be a topos, A a full generating subcategory}
\deleted{of \quad , made up with 0-connected objects. We assume \quad totally 0-}
\deleted{connected or (what amounts to the same) that the product of any two}
\deleted{objects in A is 0-connected, and moreover that any``non-empty'' object}
\deleted{of \quad has a section (again it is enough to check this for the objects}
\deleted{in A). Then the canonical homotopy structure h of \quad (defined in}
\deleted{terms of intervals which are 0-connected as a generating family of}
\deleted{homotopy intervals) is equal to the homotopy structure \hand{\unsure{h}} generated}

\margin{31}
